\documentclass[11pt]{article}
\usepackage[margin=0.9in,letterpaper]{geometry}
\usepackage{enumitem}
\usepackage{amsmath}
\usepackage{parskip}
\usepackage[colorlinks=true,linkcolor=black]{hyperref}
\setlist[enumerate]{itemsep=7pt, topsep=5pt}
\pagestyle{plain}
\begin{document}

\begin{center}
{\Large\textbf{Stat 220: A Map of Models}}\\[3pt]
{\large Midterm practice questions, with answers}\\[6pt]
\end{center}

Practice on what separates a probability model from an algorithmic one, what each lets you ask,
and how the job you were given narrows the choice. No computer needed.

\vspace{8pt}
\hrule
\vspace{10pt}

\begin{enumerate}[itemsep=14pt]
\item What makes something a probability model rather than an algorithmic one?

\begin{enumerate}[label=\Alph*., itemsep=1pt, topsep=3pt]
  \item It uses probability somewhere in the fitting procedure.
  \item It writes down a distribution for the outcome, so the data has a stated origin.
  \item It produces predictions between 0 and 1.
  \item It was fit with maximum likelihood rather than least squares.
\end{enumerate}

\vspace{2pt}
\textbf{Answer: B.} A probability model says where the data came from: `mpg ~ Normal(b0 + b1*weight, sigma^2)`. That single sentence is what produces a likelihood, and the likelihood is what standard errors, p-values, confidence intervals and AIC are all built out of. Option D describes a fitting method, which is downstream of the distribution, not the thing that defines it.
\vspace{4pt}

\item Your colleague asks for the $p$-value on the most important variable in a random forest. What do you tell them?

\begin{enumerate}[label=\Alph*., itemsep=1pt, topsep=3pt]
  \item It is in the model output, you just have to request it.
  \item You can get one by bootstrapping the forest 1,000 times.
  \item There isn't one. A forest states no distribution, so there is no likelihood and no $p$-value to report.
  \item The $p$-value is 1 divided by the number of trees.
\end{enumerate}

\vspace{2pt}
\textbf{Answer: C.} The quantity does not exist rather than being hidden. You can certainly resample to see how much an importance score wobbles, which is worth doing, but that is a stability check and not a test of a null hypothesis. There is no null hypothesis here to test.
\vspace{4pt}

\item You want to compare a linear regression against a gradient boosted model. Which comparison works?

\begin{enumerate}[label=\Alph*., itemsep=1pt, topsep=3pt]
  \item AIC, since it is designed for comparing models.
  \item Adjusted $R^2$ on the training data.
  \item Cross-validated error, using the same folds for both.
  \item Neither can be compared, they are different kinds of model.
\end{enumerate}

\vspace{2pt}
\textbf{Answer: C.} AIC needs a likelihood, so boosting is not eligible. Out-of-sample error works for anything that can produce a prediction, which is why it is the general-purpose answer. The catch is that both models have to face identical folds, otherwise you are partly measuring which model got the easier split.
\vspace{4pt}

\item Which of these is a hyperparameter?

\begin{enumerate}[label=\Alph*., itemsep=1pt, topsep=3pt]
  \item The slope on square footage in a rent model.
  \item The maximum depth of a decision tree.
  \item The residual standard deviation in a linear regression.
  \item The predicted value for a new house.
\end{enumerate}

\vspace{2pt}
\textbf{Answer: B.} A parameter comes out of fitting. A hyperparameter you set before fitting, and no amount of fitting will discover it. Depth, the lasso penalty, and $k$ in nearest neighbours are all choices you make before the data gets a say. How to make them without fooling yourself is Unit 4.
\vspace{4pt}

\item A hospital dataset has a yes/no outcome. Which of these does that fact rule out?

\begin{enumerate}[label=\Alph*., itemsep=1pt, topsep=3pt]
  \item Random forests, since they are built for numeric outcomes.
  \item Linear regression, because it cannot keep its predictions between 0 and 1.
  \item Gradient boosting, which only works on counts.
  \item Nothing. Every model can handle any outcome type.
\end{enumerate}

\vspace{2pt}
\textbf{Answer: B.} The type of the outcome rules out probability models, one at a time: linear regression genuinely cannot take a yes/no, and Poisson genuinely cannot take a continuous number. It rules out nothing on the algorithmic side, because trees, forests and boosting all have a classification mode. So when you decide against a forest, the reason is never the outcome. It is that somebody needs to read a number out of the model.
\vspace{4pt}

\item You run a lasso on 200 candidate variables, 12 survive, and you report their coefficients with $p$-values. What is wrong?

\begin{enumerate}[label=\Alph*., itemsep=1pt, topsep=3pt]
  \item Nothing, the lasso reports valid standard errors.
  \item The lasso shrinks every coefficient toward zero on purpose, so the printed values are not describing what they claim, and the variables were chosen by looking at the data.
  \item You should have used ridge instead, which does report valid $p$-values.
  \item The problem is only that 12 is too few survivors to be meaningful.
\end{enumerate}

\vspace{2pt}
\textbf{Answer: B.} Two things are broken at once. The penalty deliberately biases the coefficients, and the selection step means these 12 are the winners of a 200-way search, so their inference is computed as if the model had been chosen in advance when it was not. Use the lasso to screen, computed as if the model had been chosen in advance when it was not. Treat what survives as a shortlist worth looking into, and say that a search produced it.
\vspace{4pt}

\item A regulator will review every loan rejection your model produces. What does that fact alone tell you about the model choice?

\begin{enumerate}[label=\Alph*., itemsep=1pt, topsep=3pt]
  \item Nothing, accuracy is what matters and you should pick the most accurate model.
  \item You need a probability model, because somebody will read a coefficient and you must be able to defend it.
  \item You need the largest possible training set.
  \item You should use a neural network, since they handle complex rules.
\end{enumerate}

\vspace{2pt}
\textbf{Answer: B.} The audience decided this before you looked at a single row. A boosted model might predict default better and still be unusable, because 'the algorithm said so' is not a defense. This is the fifth job from the unit, deployment, and it constrains the choice more tightly than accuracy does.
\vspace{4pt}

\item You have 300 rows and 80 candidate columns. What do you rule out immediately?

\begin{enumerate}[label=\Alph*., itemsep=1pt, topsep=3pt]
  \item Linear regression, because 80 columns is too many for it.
  \item Flexible models like forests and boosting, because there is not enough data to pay for that much freedom.
  \item Nothing, all models work at any sample size.
  \item Cross-validation, because the folds would be too small.
\end{enumerate}

\vspace{2pt}
\textbf{Answer: B.} Flexibility is bought with sample size. With 300 rows and 80 columns a flexible model will fit the data beautifully and predict new cases badly, because most of what it is fitting is noise. Regularized regression is the sensible move here, used to cut 80 columns down to a handful chosen for a reason.
\vspace{4pt}

\item Your forest beats your regression by five points of cross-validated error. What should you do next?

\begin{enumerate}[label=\Alph*., itemsep=1pt, topsep=3pt]
  \item Ship the forest, it won.
  \item Ship the regression, simpler is always better.
  \item Find out what the forest is using that the regression is not, then decide.
  \item Average the two models together.
\end{enumerate}

\vspace{2pt}
\textbf{Answer: C.} A gap that size is information rather than a verdict. Usually the forest has found an interaction or a threshold you can name, and once you put it in the regression the gap closes. Then you ship the model you can explain. If the gap does not close, you have learned that the structure is genuinely complicated, which is also worth knowing.
\vspace{4pt}

\item A depth-10 tree has a training error of 0.44 and a cross-validated error of 13.19. A straight line has 11.59 and 12.23. Which is the better model?

\begin{enumerate}[label=\Alph*., itemsep=1pt, topsep=3pt]
  \item The tree, since 0.44 is by far the lowest number on the table.
  \item The line, because the tree only looks good on the rows it was fitted to.
  \item Neither, the two cannot be compared.
  \item The tree, because it has more flexibility available.
\end{enumerate}

\vspace{2pt}
\textbf{Answer: B.} Training error measures how well a model reproduces data it has already seen, and a flexible model can drive that near zero by memorizing. The honest column is the other one, and there the tree is worse than a straight line. Notice also that the two models sit on the same rows and are scored the same way, which is what makes the comparison mean anything.
\vspace{4pt}

\item A model predicting daily revenue reports a cross-validated MSE of 140. What should you tell the owner?

\begin{enumerate}[label=\Alph*., itemsep=1pt, topsep=3pt]
  \item The model is off by about \$140 on a typical day.
  \item The model is off by about \$12 on a typical day, since $\sqrt{140} \approx 12$.
  \item The model explains 140\% of the variation.
  \item MSE cannot be translated into dollars.
\end{enumerate}

\vspace{2pt}
\textbf{Answer: B.} MSE is in squared units, so 140 is 140 squared dollars, which means nothing to anyone. Taking the square root puts it back into dollars, and about \$12 a day is a sentence the owner can actually use. Quote the RMSE for that reason, and keep the MSE for comparing models to each other.
\vspace{4pt}

\item Which of these is decided by fitting the model, rather than by you?

\begin{enumerate}[label=\Alph*., itemsep=1pt, topsep=3pt]
  \item How deep to let a tree grow.
  \item How many folds to use in cross-validation.
  \item Where a tree puts its split points.
  \item How large a penalty to put on the coefficients.
\end{enumerate}

\vspace{2pt}
\textbf{Answer: C.} Fitting is the part the data decides: the split points, the slopes, the weights. Depth, the number of folds and the size of a penalty are all settings you choose beforehand, and no amount of fitting will choose them for you. Unit 4 covers how to make those choices without fooling yourself.
\vspace{4pt}

\end{enumerate}

\end{document}
